% Fichier généré automatiquement par tools/generate_corrections.py.
% Source : CIN/CIN-02-VitesseAcceleration/05_RT_02/05_RT_02.tex
% Statut : complet (4/4 question(s) renseignée(s), 0 reprise(s)).
\begin{corrigebox}[Corrigé extrait de la banque]
\CorrigeQuestion{1}
$\vectv{B}{2}{0} = \deriv{\vect{AB}}{\rep{0}}$ $=\deriv{\lambda(t)\vect{i_1}}{\rep{0}}$
$=\dot{\lambda}(t)\vect{i_1}+\lambda(t)\dot{\theta}(t)\vect{j_1}$.
\CorrigeQuestion{2}
$\vectv{B}{2}{0}=\vectv{B}{2}{1}+\vectv{B}{1}{0}$.

$\forall P$, $\vectv{P}{2}{1}=\dot{\lambda}(t)\vi{1}$.

Par ailleurs $\babarv{B}{A}{1}{0}=-\lambda(t)\vi{1}\wedge\dot{\theta}(t)\vk{0}$
$=\lambda(t)\dot{\theta}(t)\vj{1}$.

Au final, $\vectv{B}{2}{0}=\dot{\lambda}(t)\vi{1}+\lambda(t)\dot{\theta}(t)\vj{1}$.
\CorrigeQuestion{3}
$\torseurcin{V}{2}{0} = \torseurl{\dot{\theta}(t)\vk{0}}{\dot{\lambda}(t)\vi{1}+\lambda(t)\dot{\theta}(t)\vj{1}}{B}$.
\CorrigeQuestion{4}
$\vectg{B}{2}{0} = \deriv{\vectv{B}{2}{0}}{\rep{0}}$

$ = \ddot{\lambda}(t)\vi{1}+  \dot{\lambda}(t)\dot{\theta}\vj{1}+\dot{\lambda}(t)\dot{\theta}\vj{1}-\lambda(t)\dot{\theta}(t)^2\vi{1}$

$ = \left( \ddot{\lambda}(t)-\lambda(t)\dot{\theta}(t)^2\right)\vi{1}  +  2 \dot{\lambda}(t)\dot{\theta}(t) \vj{1}$.
\end{corrigebox}
